\subsection{精力很低的时候，应该怎样行动？}

如果你的意志能量很低，不要要求自己完成一次英雄式重启。
先把目标状态缩小到仍然有机会击败框架惯性的程度，并选用你能稳定进入的最小决策窗口。
业务目标不是证明自己很强，而是保住运动方向。
在实践上，这意味着：选一个足够便宜、足够可见、又足以为下一步建立落脚点的具体动作。

\begin{enumerate}[nosep]
  \item 用一句话命名目标状态，但把它缩小到最小的有用版本。如果“写完报告”太重，就改成“打开文件并写前三行”。
  \item 在尝试前先去掉一个明显障碍。移走手机、打开文档、把书放上桌，或提前准备工具，让当前状态失去一部分支撑。
  \item 选那个最可见、最不容易被误解的第一动作。一个弱起步也比模糊意图强。
  \item 立刻进入决策窗口。低精力意味着拖延会更贵，因为旧状态会很快重新夺回控制。
  \item 如果需要，就在第一个稳定落脚点后停下。目标是创建一条新轨迹，而不是从疲惫系统里榨出最大产出。
  \item 下一次尽量重复同样的进入模式，让这次转变越来越便宜。比起一次戏剧性的努力，小的历史累积更重要。
\end{enumerate}

\begin{remark}[低精力规则]
当精力很低时，正确问题不是“我怎么把所有事都做完”，而是“今天什么样的最小目标状态还能赢过当前状态？”
这个问题能把读者留在可控制的范围内，在这里，障碍降低和时机选择比强度更重要。
\end{remark}

\begin{example}[疲惫的夜晚]
一个读者下班回家后非常疲惫，但又想学习。
他没有安排完整的两小时会话，而是把目标状态改成“读一页并写一个问题”。
晚饭前先把书放到桌上，手机静音移开，并在饭后最安静的第一分钟里开始。
即使整段会话最后只持续了十分钟，这次转变依然是成功的，因为新框架已经在旧框架最弱的时候被启动了。
\end{example}